% Part: history 
% Chapter: biographies 
% Section: julia-robinson
\documentclass[../../../include/open-logic-section]{subfiles}

\begin{document}

\olfileid{his}{bio}{rob}

\olsection{જુલિયા રોબિન્સન}

\olphoto{robinson-julia}{જુલિયા રોબિન્સન}

જુલિયા બોમન રોબિન્સન અમેરિકન ગણિતશાસ્ત્રી
હતાં. તેઓ મુખ્યત્વે નિર્ણય સમસ્યાઓ પરના
કાર્ય માટે અને ખાસ કરીને હિલ્બર્ટની
દસમી સમસ્યાના ઉકેલમાં પોતાના પ્રદાન
માટે જાણીતા છે. રોબિન્સનનો જન્મ
8 ડિસેમ્બર 1919ના રોજ મિઝૂરીના
સેન્ટ લૂઇસમાં થયો હતો. બાળપણથી જ
સંખ્યાઓ પ્રત્યે આકર્ષણ હતું એવું
તેઓ યાદ કરતાં \citep[4]{Reid1986}.
નવ વર્ષની ઉંમરે તેમને લાલ તાવ
થયો અને ત્યાર પછી સંધિવા તાવના
વારંવાર હુમલા થયા. તેથી ઘણો
સમય પથારીમાં વિતાવવો પડ્યો અને
તેઓ અભ્યાસમાં પાછળ પડી ગયાં.
ખાનગી શિક્ષકોની મદદથી અભ્યાસમાં
બીજાઓની બરાબરી કરી શક્યાં,
પરંતુ બીમારીની શારીરિક અસરો
તેમના જીવનમાં લાંબા સમય
સુધી રહી.

બાળપણની મુશ્કેલીઓ છતાં
રોબિન્સને ગણિત અને વિજ્ઞાનનાં
ઘણાં પારિતોષિકો સાથે માધ્યમિક
શાળાનો અભ્યાસ પૂરો કર્યો.
તેમણે સાન ડિએગો સ્ટેટ કૉલેજમાં
ઉચ્ચ અભ્યાસ શરૂ કર્યો અને
છેલ્લા વર્ષમાં યુનિવર્સિટી ઑફ
કૅલિફોર્નિયા, બર્કલીમાં ગયાં.
ત્યાં ગણિતશાસ્ત્રી રાફેલ
રોબિન્સનનો તેમના પર પ્રભાવ
પડ્યો. બંને સારા મિત્ર બન્યાં
અને 1941માં લગ્ન કર્યાં.
અધ્યાપકના જીવનસાથી હોવાથી
જુલિયાને બર્કલીના ગણિત
વિભાગમાં ભણાવવા દેવામાં
આવતું ન હતું. ગણિતના
વર્ગોમાં નિયમિત પ્રવેશ વિના
હાજરી આપવાનું ચાલુ રાખ્યું
હોવા છતાં તેઓ યુનિવર્સિટી
છોડી પરિવાર શરૂ કરવા માગતાં
હતાં. જોકે લગ્ન પછી થોડા
જ સમયમાં તેમને ન્યુમોનિયા
થયો. બાળપણના સંધિવા
તાવને કારણે તેમના હૃદય
પર નોંધપાત્ર પ્રમાણમાં
ડાઘવાળી પેશી બની હતી એવું તેમને
કહેવામાં આવ્યું. આ
અસરની ગંભીરતાને કારણે
ડૉક્ટરે તેઓ ચાલીસ વર્ષથી
વધુ નહીં જીવે એવી આગાહી
કરી અને સંતાન ન કરવાની
સલાહ આપી \citep[13]{Reid1986}.

રોબિન્સન લાંબા સમય સુધી
અવસાદમાં રહ્યાં, પણ આખરે
ગણિતનો અભ્યાસ ચાલુ રાખવાનું
નક્કી કર્યું. તેઓ બર્કલી
પાછાં ગયાં અને 1948માં
આલ્ફ્રેડ ટાર્સ્કીના માર્ગદર્શન
હેઠળ પીએચ.ડી. પૂરી કરી.
ટાર્સ્કીએ બતાવ્યું હતું કે
વાસ્તવિક સંખ્યાઓનો પ્રથમ-ક્રમ
સિદ્ધાંત નિર્ણેય છે, જ્યારે
ગેડલના કાર્યમાંથી ફલિત થતું
હતું કે પ્રાકૃતિક સંખ્યાઓનો
પ્રથમ-ક્રમ સિદ્ધાંત અનિર્ણેય
છે. સંમેય સંખ્યાઓનો પ્રથમ-ક્રમ
સિદ્ધાંત નિર્ણેય છે કે નહીં
તે મોટો ખુલ્લો પ્રશ્ન હતો.
પોતાના નિબંધમાં
\citeyearpar{Robinson1949}
રોબિન્સને સાબિત કર્યું કે
તે અનિર્ણેય છે.

નિર્ણય સમસ્યાઓમાં રસ હોવાથી
રોબિન્સને પછી હિલ્બર્ટની
દસમી સમસ્યાનો ઉકેલ શોધવાનો
પ્રયત્ન કર્યો. ડેવિડ હિલ્બર્ટે
1900માં રજૂ કરેલી 23 પ્રસિદ્ધ
ગાણિતિક સમસ્યાઓમાંની આ
એક હતી. દસમી સમસ્યા પૂછે
છે કે પૂર્ણાંક ગુણાંકોવાળા
બહુપદી સમીકરણને, જેમ કે
$3x^2 - 2y +3 = 0$ને,
પૂર્ણાંકોમાં ઉકેલ છે કે
નહીં તે સાન્ત સમયમાં
નક્કી કરનારી કલનવિધિ
છે કે નહીં. આવા પ્રશ્નો
\emph{ડાયોફૅન્ટાઇન સમસ્યાઓ}
કહેવાય છે. પ્રારંભિક
સફળતાઓ પછી રોબિન્સન
આ સમસ્યા પર કામ કરતા
માર્ટિન ડેવિસ અને
હિલરી પુટનામ સાથે
જોડાયાં. તેમણે બતાવ્યું
કે ઘાતીય ડાયોફૅન્ટાઇન
સમસ્યાઓ, જેમાં અજ્ઞાત
રાશિઓ ઘાતાંક તરીકે પણ
આવી શકે છે, અનિર્ણેય
છે. તેમણે એ પણ
બતાવ્યું કે પછીથી
“J.R.” કહેવાયેલું
એક ચોક્કસ અનુમાન
હિલ્બર્ટની દસમી
સમસ્યા અનિર્ણેય
છે એવું ફલિત કરે
છે \citep{DavisPutnamRobinson1961}.
રોબિન્સન 1960ના
દાયકામાં આ સમસ્યા
પર કામ કરતાં રહ્યાં.
1970માં યુવાન રશિયન
ગણિતશાસ્ત્રી યુરી
માતિયાસેવિચે છેવટે
J.R. અનુમાન સાબિત
કર્યું. સંયુક્ત પરિણામ
આજે માતિયાસેવિચ--રોબિન્સન--
ડેવિસ--પુટનામ પ્રમેય,
અથવા ટૂંકમાં MRDP પ્રમેય,
કહેવાય છે. માતિયાસેવિચ
અને રોબિન્સન મિત્ર
બન્યાં અને અનેક લેખો
પર સાથે કામ કર્યું.
માતિયાસેવિચને લખેલા
પત્રમાં રોબિન્સને
એક વાર લખ્યું હતું:
“ખરેખર મને બહુ આનંદ
છે કે સાથે કામ કરીને
(હજારો માઇલ દૂર
હોવા છતાં) આપણે
બંનેમાંથી કોઈ એકલા
કરી શકે તેના કરતાં
દેખીતી રીતે વધુ
પ્રગતિ કરી રહ્યા છીએ”
\citep[45]{Matijasevich1992}.

રોબિન્સન અમેરિકન
ગણિતીય મંડળનાં
પ્રથમ મહિલા પ્રમુખ
હતાં અને રાષ્ટ્રીય
વિજ્ઞાન અકાદમીમાં
ચૂંટાયેલાં પણ પ્રથમ
સ્ત્રી હતાં.
લ્યુકેમિયાનું નિદાન
થયા પછી 30 જુલાઈ
1985ના રોજ 65
વર્ષની ઉંમરે
તેમનું અવસાન થયું.

\begin{reading}
રોબિન્સનના ગાણિતિક
લેખો તેમના
\textit{Collected
  Works}માં ઉપલબ્ધ છે
\citep{Robinson1996}.
તેમાં રાષ્ટ્રીય વિજ્ઞાન
અકાદમી માટે લખાયેલ
તેમનું જીવનચરિત્રાત્મક
સ્મૃતિલેખ પણ ફરી
છાપવામાં આવ્યો છે
\citep{Feferman1994}.
રોબિન્સનનાં મોટાં
બહેન કૉન્સ્ટન્સ રીડે
મુલાકાતોના આધારે
“Autobiography of Julia”
પ્રકાશિત કર્યું
\citep{Reid1986}
અને સંપૂર્ણ
સ્મૃતિલેખ પણ લખ્યો
\citep{Reid1996}.
જૉર્જ સિસસેરીએ
રોબિન્સન અને
હિલ્બર્ટની દસમી
સમસ્યા વિશે ટૂંકો
દસ્તાવેજી ચલચિત્ર
બનાવ્યું
\citep{Csicsery2016}.
માતિયાસેવિચ સાથે
રોબિન્સનના સહકાર
અને તેમના કાર્ય પર
રોબિન્સનની અસર
વિશેના ટૂંકા
સ્મૃતિલેખ માટે
\citep{Matijasevich1992}
જુઓ.
\end{reading}

\end{document}
